\documentclass[aspectratio=169]{beamer}

\usetheme{Madrid}
\usecolortheme{default}
\usefonttheme{professionalfonts}

\usepackage[T1]{fontenc}
\usepackage{lmodern}
\usepackage{amsmath,amssymb}
\usepackage{bm}
\usepackage{booktabs}
\usepackage{xcolor}

\title[PQC Hard Problem Colloquium]{Orientation Materials for the PQC Hard Problem Colloquium}
\author{Your Name}
\institute{Your Institution / Research Group}
\date{\today}

\begin{document}
	
	\begin{frame}
		\titlepage
	\end{frame}
	
	\begin{frame}{Purpose of This Colloquium}
		\begin{itemize}
			\item This colloquium is about the \textbf{hard problems} that underlie modern public-key cryptography and post-quantum cryptography.
			\item The goal is not only to know \emph{which schemes exist}, but to understand:
			\begin{itemize}
				\item what mathematical problem each scheme relies on,
				\item what the best known attacks are,
				\item how classical and quantum complexity differ,
				\item and why some families remain viable while others have failed.
			\end{itemize}
			\item In other words, the central object of study is the \textbf{security assumption itself}.
		\end{itemize}
	\end{frame}
	
	\begin{frame}{Guiding Question}
		\begin{block}{Core Question}
			Which computational problems remain hard enough to support secure cryptography in the presence of quantum computers?
		\end{block}
		
		\vspace{0.5em}
		This question naturally divides our discussion into three layers:
		\begin{enumerate}
			\item the \textbf{mathematical hard problem},
			\item the \textbf{best known attacks},
			\item the \textbf{cryptosystems} built on top of that assumption.
		\end{enumerate}
	\end{frame}
	
	\begin{frame}{Why Hard Problems Matter}
		\begin{itemize}
			\item A cryptosystem is only as trustworthy as the computational problem beneath it.
			\item Cryptographic history shows that assumptions can fail in two very different ways:
			\begin{itemize}
				\item through a major algorithmic breakthrough,
				\item or through hidden algebraic structure that makes a scheme easier than expected.
			\end{itemize}
			\item Post-quantum cryptography therefore requires a careful comparison of:
			\begin{itemize}
				\item \textbf{classical hardness},
				\item \textbf{quantum hardness},
				\item \textbf{practical attack cost},
				\item and \textbf{structural fragility}.
			\end{itemize}
		\end{itemize}
	\end{frame}
	
	\section{From Classical Assumptions to the Post-Quantum Setting}
	
	\begin{frame}{The Classical Baseline}
		Before post-quantum cryptography, public-key cryptography was dominated by a small number of number-theoretic assumptions:
		\begin{itemize}
			\item \textbf{Integer Factorization} $\rightarrow$ RSA, Rabin
			\item \textbf{Discrete Logarithm in Finite Fields} $\rightarrow$ Diffie--Hellman, DSA, ElGamal
			\item \textbf{Elliptic Curve Discrete Logarithm} $\rightarrow$ ECDH, ECDSA, EdDSA
		\end{itemize}
		
		\vspace{0.5em}
		These problems were historically attractive because they enabled elegant and efficient cryptographic constructions.
	\end{frame}
	
	\begin{frame}{Why the Classical Baseline Fails in the Quantum Model}
		\begin{itemize}
			\item Classically, factorization and finite-field discrete logarithm admit only \textbf{sub-exponential} attacks in their best known forms.
			\item Elliptic-curve discrete logarithm is even stronger classically, with generic attacks of roughly \textbf{$O(\sqrt{n})$}.
			\item However, in the quantum setting, all three become vulnerable to \textbf{Shor's algorithm}.
		\end{itemize}
		
		\vspace{0.75em}
		\begin{block}{Consequent Shift}
			The PQC transition is fundamentally the replacement of Shor-vulnerable assumptions with new assumptions for which no efficient quantum attacks are known.
		\end{block}
	\end{frame}
	
	\begin{frame}{A Historical Security Lesson}
		The classical world teaches an important lesson for PQC:
		\begin{itemize}
			\item security is not only about asymptotic hardness,
			\item it is also about the existence of \textbf{unexpected algorithmic structure},
			\item and about whether the problem remains hard under the exact distribution used by the scheme.
		\end{itemize}
		
		\vspace{0.5em}
		This is why the colloquium must focus on \textbf{hard problems first}, not only on named standards or implementations.
	\end{frame}
	
	\section{Main Families of Post-Quantum Hard Problems}
	
	\begin{frame}{The Main PQC Families}
		The major families discussed in post-quantum cryptography are:
		\begin{enumerate}
			\item \textbf{Lattice-based problems}
			\item \textbf{Code-based problems}
			\item \textbf{Hash-based assumptions}
			\item \textbf{Multivariate polynomial problems}
			\item \textbf{Isogeny-based problems}
		\end{enumerate}
		
		\vspace{0.5em}
		Each family has a different profile with respect to:
		\begin{itemize}
			\item attack methodology,
			\item concrete efficiency,
			\item reduction theory,
			\item and quantum resistance.
		\end{itemize}
	\end{frame}
	
	\subsection{Lattice-Based Problems}
	
	\begin{frame}{Lattice Problems}
		The lattice family includes problems such as:
		\begin{itemize}
			\item \textbf{LWE} (Learning With Errors),
			\item \textbf{SIS} (Short Integer Solution),
			\item \textbf{SVP} (Shortest Vector Problem),
			\item \textbf{CVP} (Closest Vector Problem).
		\end{itemize}
		
		\vspace{0.5em}
		These assumptions support schemes such as:
		\begin{itemize}
			\item Kyber / ML-KEM,
			\item Dilithium / ML-DSA,
			\item NTRU,
			\item fully homomorphic encryption constructions.
		\end{itemize}
	\end{frame}
	
	\begin{frame}{Why Lattices Matter}
		Lattice-based cryptography is central because:
		\begin{itemize}
			\item it is currently the most influential practical PQC family,
			\item it supports both encryption/KEM and signature constructions,
			\item and it has strong connections between average-case and worst-case hardness in several settings.
		\end{itemize}
		
		\vspace{0.5em}
		From the attack perspective, the dominant language is:
		\begin{itemize}
			\item \textbf{lattice reduction} (LLL, BKZ),
			\item \textbf{sieving},
			\item and careful \textbf{parameter estimation}.
		\end{itemize}
	\end{frame}
	
	\begin{frame}{Lattice Attack Perspective}
		According to the table, the best known attacks remain broadly \textbf{exponential} in the relevant dimension parameter.
		
		\vspace{0.5em}
		This is the key orientation point:
		\begin{itemize}
			\item there is no Shor-style polynomial-time quantum algorithm known for the main lattice assumptions,
			\item but concrete security depends heavily on accurate cost models for BKZ and related reductions,
			\item so lattice security is a field where asymptotics and concrete cryptanalysis must be studied together.
		\end{itemize}
	\end{frame}
	
	\subsection{Code-Based Problems}
	
	\begin{frame}{Code-Based Problems}
		Code-based cryptography is centered on problems such as:
		\begin{itemize}
			\item \textbf{syndrome decoding},
			\item more generally, the difficulty of decoding random linear codes.
		\end{itemize}
		
		\vspace{0.5em}
		Representative systems include:
		\begin{itemize}
			\item McEliece,
			\item HQC,
			\item BIKE.
		\end{itemize}
	\end{frame}
	
	\begin{frame}{Code-Based Attack Perspective}
		The dominant attack family is \textbf{Information Set Decoding (ISD)} and its variants.
		
		\vspace{0.5em}
		The main qualitative lesson is:
		\begin{itemize}
			\item the best known attacks remain \textbf{exponential},
			\item quantum speedups appear much more limited than in factorization or DLP,
			\item and Grover-type improvements generally give only a \textbf{quadratic speedup}, rather than a full structural collapse.
		\end{itemize}
		
		\vspace{0.5em}
		This makes code-based cryptography one of the most conservative PQC directions.
	\end{frame}
	
	\subsection{Hash-Based Assumptions}
	
	\begin{frame}{Hash-Based Assumptions}
		Hash-based cryptography relies on familiar assumptions:
		\begin{itemize}
			\item \textbf{collision resistance},
			\item \textbf{preimage resistance},
			\item \textbf{second-preimage resistance}.
		\end{itemize}
		
		\vspace{0.5em}
		Representative systems include:
		\begin{itemize}
			\item SPHINCS+,
			\item XMSS,
			\item LMS.
		\end{itemize}
	\end{frame}
	
	\begin{frame}{Hash-Based Security Perspective}
		The attack picture is comparatively clean:
		\begin{itemize}
			\item collisions face the \textbf{birthday bound},
			\item preimage search is affected by \textbf{Grover's algorithm},
			\item but there is no analogous Shor-style collapse.
		\end{itemize}
		
		\vspace{0.5em}
		Therefore, hash-based signatures are often viewed as:
		\begin{itemize}
			\item conceptually conservative,
			\item relatively robust at the assumption level,
			\item though sometimes less attractive in size or efficiency.
		\end{itemize}
	\end{frame}
	
	\subsection{Multivariate Problems}
	
	\begin{frame}{Multivariate Polynomial Problems}
		Multivariate cryptography is built on the difficulty of solving systems of quadratic equations over finite fields.
		
		\vspace{0.5em}
		This is often referred to as the \textbf{MQ problem}.
		
		\vspace{0.5em}
		Representative systems include:
		\begin{itemize}
			\item UOV,
			\item MAYO,
			\item historically, Rainbow.
		\end{itemize}
	\end{frame}
	
	\begin{frame}{Multivariate Attack Perspective}
		The table highlights the relevant attack language:
		\begin{itemize}
			\item \textbf{MinRank}-type attacks,
			\item \textbf{Kipnis--Shamir}-type methods,
			\item \textbf{XL / Gr\"obner basis} approaches,
			\item and scheme-specific structural attacks.
		\end{itemize}
		
		\vspace{0.5em}
		The critical lesson is that multivariate systems may appear hard in the abstract, yet particular constructions can fail because of exploitable algebraic structure.
	\end{frame}
	
	\begin{frame}{Structural Fragility: The Rainbow Lesson}
		Rainbow is an important case study.
		\begin{itemize}
			\item At the high level, MQ remains a hard class of problems.
			\item But Rainbow itself was broken through a much more specific structural attack.
			\item This teaches that \textbf{scheme hardness} is not identical to \textbf{family hardness}.
		\end{itemize}
		
		\vspace{0.5em}
		This distinction will be a recurring theme throughout the colloquium.
	\end{frame}
	
	\subsection{Isogeny-Based Problems}
	
	\begin{frame}{Isogeny Problems}
		Isogeny-based cryptography emerged as a highly elegant and compact direction in PQC.
		
		\vspace{0.5em}
		Systems discussed in this space include:
		\begin{itemize}
			\item SQISign,
			\item CSIDH,
			\item and historically SIKE / SIDH.
		\end{itemize}
	\end{frame}
	
	\begin{frame}{The Isogeny Cautionary Story}
		The table captures one of the most important cautionary episodes in modern PQC:
		\begin{itemize}
			\item \textbf{SIKE is broken}.
			\item The Castryck--Decru attack showed that a scheme can collapse completely even when the surrounding research area still appears mathematically rich.
		\end{itemize}
		
		\vspace{0.5em}
		More broadly:
		\begin{itemize}
			\item some isogeny problems still resist known attacks,
			\item but the family must now be approached with increased caution,
			\item especially when hidden structure may create special-purpose vulnerabilities.
		\end{itemize}
	\end{frame}
	
	\section{Cross-Family Comparison}
	
	\begin{frame}{A Comparative Reading of the Hard Problem Landscape}
		Based on the table, the families differ in a striking way:
		
		\vspace{0.5em}
		\begin{itemize}
			\item \textbf{Factorization / DLP / ECDLP}: strong classically, but broken by polynomial-time quantum algorithms.
			\item \textbf{Lattice / Code-based}: still believed hard; attacks remain exponential or subexponential only in restricted senses.
			\item \textbf{Hash-based}: quantum attacks are limited and well understood.
			\item \textbf{Multivariate / Isogeny}: mathematically interesting, but some constructions have shown serious fragility.
		\end{itemize}
	\end{frame}
	
	\begin{frame}{Three Evaluation Questions for Every Scheme}
		For each cryptographic construction studied in this colloquium, we should ask:
		
		\begin{enumerate}
			\item \textbf{What exact hard problem is being assumed?}
			\item \textbf{What are the best known attacks, classically and quantumly?}
			\item \textbf{Does the scheme introduce exploitable structure beyond the underlying problem?}
		\end{enumerate}
		
		\vspace{0.75em}
		This framework helps separate genuine hardness from merely optimistic design.
	\end{frame}
	
	\begin{frame}{What the Table Really Teaches}
		The information in the table should not be read as a static catalog.
		
		\vspace{0.5em}
		It should be read as a set of research attitudes:
		\begin{itemize}
			\item compare assumptions by their \textbf{attack surface},
			\item distinguish \textbf{family-level hardness} from \textbf{scheme-level security},
			\item treat broken constructions as lessons, not merely failures,
			\item and always ask whether the best known attack is generic, structural, classical, or quantum.
		\end{itemize}
	\end{frame}
	
	\section{Orientation for the Colloquium}
	
	\begin{frame}{How This Colloquium Can Be Structured}
		A natural progression for the colloquium is:
		\begin{enumerate}
			\item Classical hard problems and the quantum threat
			\item Lattice assumptions and reduction-based security
			\item Code-based cryptography and ISD
			\item Hash-based signatures as a conservative baseline
			\item Multivariate cryptography and structural attacks
			\item Isogeny cryptography after SIKE
			\item Comparative security evaluation across families
		\end{enumerate}
	\end{frame}
	
	\begin{frame}{Expected Outcomes}
		By the end of the colloquium, participants should be able to:
		\begin{itemize}
			\item explain the difference between classical and post-quantum hardness assumptions,
			\item describe the principal attack models against each family,
			\item identify why some assumptions remain credible and others do not,
			\item and evaluate a PQC scheme in terms of both mathematics and cryptanalysis.
		\end{itemize}
	\end{frame}
	
	\begin{frame}{Closing Perspective}
		\begin{block}{Final Perspective}
			Post-quantum cryptography is not merely a search for new schemes. It is a disciplined search for mathematical problems whose hardness survives new computational models, new attack paradigms, and deeper structural analysis.
		\end{block}
		
		\vspace{0.75em}
		That is the central intellectual purpose of the PQC Hard Problem Colloquium.
	\end{frame}
	
	\begin{frame}
		\centering
		{\LARGE Thank you.}
		
		\vspace{0.5em}
		{\large Discussion and questions}
	\end{frame}
	
\end{document}